\documentclass[12pt, a4paper]{article}

\usepackage[utf8]{inputenc}
\usepackage[T1]{fontenc}
\usepackage[brazil]{babel}
\usepackage{amsmath, amssymb}
\usepackage{graphicx}
\usepackage{float}
\usepackage{booktabs}
\usepackage{geometry}
\usepackage{enumitem}
\usepackage{setspace}
\usepackage{fancyhdr}
\usepackage[hyphens]{url}
\usepackage{hyperref}

\geometry{left=2.5cm, right=2.5cm, top=2.5cm, bottom=2.5cm}
\setstretch{1.3}
\hypersetup{
    colorlinks=true,
    linkcolor=blue!70!black,
    citecolor=blue!70!black,
    urlcolor=blue!70!black
}

\pagestyle{fancy}
\fancyhf{}
\rhead{Projeto 1 — Modelo Teórico Prophet}
\lhead{Séries Temporais}
\cfoot{\thepage}

\begin{document}

% ---- CAPA ----
\begin{titlepage}
\centering
\vspace*{2cm}
{\Large \textbf{UNIVERSIDADE ESTADUAL DE CAMPINAS}}\\[0.2cm]
{\large Séries Temporais}\\[3cm]
{\LARGE \textbf{Modelo Prophet para Séries Temporais:\\
Teoria, Ilustração e Discussão Crítica}}\\[2cm]
{\Large Projeto 1 — Modelo Teórico}\\[3cm]
{\large
\textbf{Nome:} Arthur Lourenço Narcizo da Silva\\[0.2cm]
\textbf{RA:} 236043\\[0.2cm]
}
{\large Campinas, \today}
\end{titlepage}

\tableofcontents
\newpage

% ============================================================
\section{Introdução}

O modelo \textbf{Prophet}, proposto por Taylor e Letham (2018),
é um modelo aditivo generalizado para séries temporais com
tendência, sazonalidade e feriados, estimado via inferência
bayesiana.

Este relatório apresenta: (i)~a formulação matemática,
(ii)~uma ilustração com o índice de produção industrial
brasileiro (IBGE, Série 21859 do BCB), e (iii)~discussão crítica.

% ============================================================
\section{Formulação Teórica}

\subsection{Decomposição}

O Prophet decompõe $y(t)$ em componentes interpretáveis:
\begin{equation}
    y(t) = g(t) + s(t) + h(t) + \varepsilon_t,
    \quad \varepsilon_t \sim \mathcal{N}(0, \sigma^2)
    \label{eq:prophet}
\end{equation}
onde $g(t)$ é a tendência, $s(t)$ a sazonalidade e $h(t)$
feriados. No modo multiplicativo:
$y(t) = g(t)(1 + s(t)) + h(t) + \varepsilon_t$.

\subsection{Tendência com \textit{changepoints}}

\begin{equation}
    g(t) = \big(k + \mathbf{a}(t)^\top\boldsymbol{\delta}\big)\,t
    + \big(m + \mathbf{a}(t)^\top\boldsymbol{\gamma}\big)
    \label{eq:trend}
\end{equation}
onde $k$ é a taxa base, $m$ o intercepto,
$a_j(t) = \mathbf{1}[t \geq s_j]$, $\delta_j$ ajusta a taxa
no \textit{changepoint} $s_j$, e $\gamma_j = -s_j\delta_j$
garante continuidade. O prior regularizador é:
\begin{equation}
    \delta_j \sim \text{Laplace}(0, \tau)
\end{equation}
com $\tau$ (\texttt{changepoint\_prior\_scale}) controlando a
flexibilidade da tendência.

\subsection{Sazonalidade (Fourier)}

\begin{equation}
    s(t) = \sum_{n=1}^{N}\!\left(
    a_n\cos\!\frac{2\pi nt}{P} + b_n\sin\!\frac{2\pi nt}{P}
    \right), \quad
    \boldsymbol{\beta} \sim \mathcal{N}(\mathbf{0}, \sigma_s^2\mathbf{I})
    \label{eq:fourier}
\end{equation}
com período $P$ e ordem $N$. O prior $\sigma_s$
(\texttt{seasonality\_prior\_scale}) controla a amplitude.

\subsection{Feriados}

$h(t) = \sum_i \kappa_i \cdot \mathbf{1}[t \in D_i]$,
com $\kappa_i \sim \mathcal{N}(0, \nu^2)$.

\subsection{Estimação}

Os parâmetros são estimados via MAP usando L-BFGS no
\textit{backend} Stan (Carpenter et al., 2017):
\begin{equation}
    \hat{\boldsymbol{\theta}} = \arg\max_{\boldsymbol{\theta}}
    \big[\log p(\mathbf{y}|\boldsymbol{\theta}) +
    \log p(\boldsymbol{\theta})\big]
\end{equation}
Intervalos de confiança são gerados por simulação da incerteza
em $\boldsymbol{\delta}$, $\boldsymbol{\beta}$ e $\varepsilon_t$.

\subsection{Comparação com ARIMA}

\begin{table}[H]
\centering
\caption{Prophet vs.\ modelos do conteúdo programático.}
\begin{tabular}{lcc}
    \toprule
    \textbf{Aspecto} & \textbf{ARIMA} & \textbf{Prophet} \\
    \midrule
    Abordagem & Box-Jenkins & GAM bayesiano \\
    Estacionariedade & Requer & Não requer \\
    Sazonalidade & $(P,D,Q)_s$ & Fourier \\
    Tendência & Diferenciação & Changepoints \\
    ICs & Analíticos & Simulação \\
    Estimação & MLE & MAP (Stan) \\
    \bottomrule
\end{tabular}
\end{table}

% ============================================================
\section{Ilustração: Produção Industrial Brasileira}

\subsection{Dados}

A série utilizada é o \textbf{Índice de Produção Industrial}
(base = 100), calculado pelo \textbf{IBGE} na Pesquisa Industrial
Mensal --- Produção Física (PIM-PF), disponibilizado pelo Banco
Central do Brasil (Série 21859 do SGS).

\begin{itemize}[nosep]
    \item \textbf{Fonte}: IBGE/PIM-PF, via BCB/SGS --- Série 21859
    \item \textbf{URL}: \nolinkurl{https://api.bcb.gov.br/dados/serie/bcdata.sgs.21859/dados?formato=csv}
    \item \textbf{Variável}: Índice de produção industrial (base = 100)
    \item \textbf{Período}: janeiro de 2002 a janeiro de 2026
    \item \textbf{Observações}: 289
    \item \textbf{Frequência}: Mensal
\end{itemize}

A série reflete o nível de atividade da indústria brasileira e
apresenta tendência, sazonalidade (queda em janeiro/fevereiro e
dezembro por férias coletivas) e choques estruturais (crise de
2008, recessão 2015--2016, COVID-19 em 2020).

\begin{figure}[H]
\centering
\includegraphics[width=\textwidth]{01_serie.png}
\caption{Produção Industrial --- Brasil (IBGE, Série 21859 BCB).}
\label{fig:serie}
\end{figure}

\begin{figure}[H]
\centering
\includegraphics[width=\textwidth]{02_stl.png}
\caption{Decomposição STL.}
\label{fig:stl}
\end{figure}

O teste ADF indica não estacionariedade em nível ($p = 0{,}17$) e
estacionariedade na primeira diferença ($p = 0{,}0001$).

\begin{figure}[H]
\centering
\includegraphics[width=0.85\textwidth]{03_boxplot.png}
\caption{Sazonalidade mensal --- queda em jan., fev.\ e dez.}
\label{fig:box}
\end{figure}

\subsection{Configuração}

\begin{table}[H]
\centering
\caption{Hiperparâmetros do Prophet.}
\begin{tabular}{ll}
    \toprule
    Tendência & Linear, changepoints automáticos (25 detectados) \\
    \texttt{changepoint\_prior\_scale} & 0.1 \\
    \texttt{changepoint\_range} & 0.85 \\
    Sazonalidade & Aditiva, Fourier $N = 8$ \\
    \texttt{seasonality\_prior\_scale} & 10.0 \\
    Feriados & Carnaval (afeta produção em fev./mar.) \\
    IC & 95\% \\
    Divisão & 80\%/20\% temporal (231/58 obs.) \\
    \bottomrule
\end{tabular}
\end{table}

\subsection{Componentes}

\begin{figure}[H]
\centering
\includegraphics[width=\textwidth]{05_componentes.png}
\caption{Componentes estimados pelo Prophet.}
\label{fig:comp}
\end{figure}

\begin{figure}[H]
\centering
\includegraphics[width=\textwidth]{06_changepoints.png}
\caption{Changepoints na tendência.}
\label{fig:cp}
\end{figure}

\subsection{Resultados}

\begin{figure}[H]
\centering
\includegraphics[width=\textwidth]{07_real_vs_previsto.png}
\caption{Real vs.\ Prophet (teste) com IC 95\%.}
\label{fig:prev}
\end{figure}

\begin{table}[H]
\centering
\caption{Métricas de desempenho.}
\label{tab:metricas}
\begin{tabular}{lr}
    \toprule
    RMSE & 4.67 \\
    MAE & 3.91 \\
    MAPE & 3.79\% \\
    $R^2$ & 0.5731 \\
    $U$ de Theil & 0.0233 \\
    Cobertura IC 95\% & 100.0\% \\
    \bottomrule
\end{tabular}
\end{table}

\subsection{Benchmarks}

\begin{figure}[H]
\centering
\includegraphics[width=0.8\textwidth]{11_benchmarks.png}
\caption{MAPE: Prophet vs.\ benchmarks.}
\label{fig:bench}
\end{figure}

\begin{table}[H]
\centering
\caption{Comparação com benchmarks.}
\label{tab:bench}
\begin{tabular}{lrrr}
    \toprule
    \textbf{Modelo} & \textbf{RMSE} & \textbf{MAE} & \textbf{MAPE} \\
    \midrule
    Prophet   & 4.67 & 3.91 & 3.79\% \\
    Naïve     & 5.62 & 4.51 & 4.46\% \\
    S-Naïve   & 2.53 & 1.81 & 1.76\% \\
    Média     & 7.15 & 6.04 & 6.01\% \\
    Drift     & 7.05 & 5.97 & 5.98\% \\
    \bottomrule
\end{tabular}
\end{table}

O Prophet superou os benchmarks Naïve, Média e Drift, porém o
Sazonal-Naïve obteve MAPE inferior (1.76\% vs.\ 3.79\%),
indicando que a sazonalidade é o componente dominante nesta série
no período recente.

\subsection{Validação cruzada}

\begin{figure}[H]
\centering
\includegraphics[width=\textwidth]{10_cv_mape.png}
\caption{MAPE por horizonte (validação cruzada temporal).}
\label{fig:cv}
\end{figure}

MAPE médio da CV: 4.88\%. RMSE médio da CV: 6.84.

\subsection{Diagnósticos dos resíduos}

\begin{figure}[H]
\centering
\includegraphics[width=\textwidth]{09_diagnosticos.png}
\caption{Diagnósticos: resíduos, histograma, QQ-Plot, ACF.}
\label{fig:diag}
\end{figure}

\begin{table}[H]
\centering
\caption{Testes dos resíduos ($\alpha = 5\%$).}
\label{tab:testes}
\begin{tabular}{llcl}
    \toprule
    \textbf{Teste} & $H_0$ & \textbf{p-valor} &
    \textbf{Decisão} \\
    \midrule
    Jarque-Bera & Normalidade & 0.7003 & Não rejeita $H_0$ (Normal) \\
    Ljung-Box (12) & Sem autocorrelação & $<0.0001$ & Rejeita $H_0$ (Autocorrelação) \\
    Teste $t$ & Média $= 0$ & $<0.0001$ & Rejeita $H_0$ (Viés) \\
    Spearman & Homocedasticidade & $<0.0001$ & Rejeita $H_0$ (Heterocedasticidade) \\
    \bottomrule
\end{tabular}
\end{table}

Os resíduos apresentam distribuição normal (Jarque-Bera, $p = 0{,}70$),
porém exibem viés sistemático positivo (média $= 3{,}47$), autocorrelação
significativa e heterocedasticidade. Esses resultados confirmam a
limitação teórica do Prophet de tratar $\varepsilon_t$ como i.i.d.

% ============================================================
\section{Discussão Crítica}

\subsection{Vantagens}
\begin{enumerate}[nosep, leftmargin=*]
    \item \textbf{Decomposição interpretável}: tendência,
    sazonalidade e feriados estimados e visualizados separadamente
    --- permite ao analista entender o que impulsiona a previsão.
    \item \textbf{Changepoints automáticos}: identificou 25 mudanças
    estruturais (crise 2008, recessão 2015, COVID) sem
    intervenção manual.
    \item \textbf{ICs nativos} por simulação bayesiana, com
    cobertura de 100\% no teste.
    \item \textbf{Não requer estacionariedade}: o teste ADF
    indicou não estacionariedade em nível ($p = 0{,}17$), mas o
    Prophet não necessita de diferenciação.
    \item \textbf{Robusto a dados faltantes}.
    \item \textbf{Feriados integrados}: Carnaval afeta a produção
    industrial e é modelado diretamente.
\end{enumerate}

\subsection{Desvantagens}
\begin{enumerate}[nosep, leftmargin=*]
    \item \textbf{Não modela autocorrelação residual}:
    $\varepsilon_t$ é tratado como i.i.d. O teste de Ljung-Box
    confirmou autocorrelação significativa ($p < 0{,}0001$).
    \item \textbf{Viés sistemático}: o teste $t$ indicou média
    dos resíduos de 3.47 ($p < 0{,}0001$), ou seja, o Prophet
    subestima sistematicamente a produção no período de teste.
    \item \textbf{Sazonalidade fixa}: coeficientes de Fourier não
    mudam no tempo --- séries com sazonalidade evolutiva são mal
    representadas.
    \item \textbf{Extrapolação linear}: tendência pode divergir
    para horizontes longos.
    \item \textbf{Escolha aditivo/multiplicativo} é manual e
    impactante.
\end{enumerate}

\subsection{Limitações nesta série}
\begin{enumerate}[nosep, leftmargin=*]
    \item A produção industrial brasileira sofreu choques extremos
    (COVID-19: queda de $\approx 25\%$ em abril/2020), que o
    Prophet trata como ruído, não como evento estruturado.
    \item A tendência industrial brasileira é irregular (crescimento
    até 2013, estagnação 2014--2019, choque 2020, recuperação),
    exigindo muitos changepoints.
    \item O $R^2 = 0{,}57$ indica que o modelo explica pouco mais
    da metade da variabilidade no teste.
    \item O Sazonal-Naïve superou o Prophet (MAPE 1.76\% vs.\
    3.79\%), sugerindo que a sazonalidade domina e a tendência
    estimada introduz mais erro do que benefício.
\end{enumerate}

\subsection{Interpretação dos diagnósticos}

Os diagnósticos dos resíduos revelam três problemas:

\begin{enumerate}[nosep, leftmargin=*]
    \item \textbf{Viés sistemático}: a média dos resíduos é 3.47
    (teste $t$, $p < 0{,}0001$), indicando que o Prophet subestima
    sistematicamente a produção industrial no período de teste
    (abr/2021 a jan/2026). Isso se deve à tendência linear do
    modelo, que projeta a queda de 2020 adiante, enquanto a
    recuperação pós-COVID elevou os índices.

    \item \textbf{Autocorrelação residual}: o teste de Ljung-Box
    rejeita $H_0$ em todos os lags ($p < 0{,}0001$), confirmando
    que o Prophet não captura a dependência temporal dos erros ---
    exatamente a limitação teórica do modelo, que trata
    $\varepsilon_t$ como i.i.d.

    \item \textbf{Heterocedasticidade}: o coeficiente de Spearman
    entre $|e_t|$ e $t$ é $\rho = 0{,}69$ ($p < 0{,}0001$),
    indicando que os erros crescem ao longo do horizonte ---
    comportamento esperado quando a tendência diverge.

    \item \textbf{Normalidade}: o único teste favorável é o
    Jarque-Bera ($p = 0{,}70$), indicando que a \textit{distribuição}
    dos resíduos é aproximadamente normal, embora sua estrutura
    temporal não seja ruído branco.
\end{enumerate}

\subsection{Como melhorar}
\begin{enumerate}[nosep, leftmargin=*]
    \item Adicionar \textbf{COVID-19 como feriado} com
    \texttt{lower\_window} e \texttt{upper\_window} para
    modelar explicitamente o choque.
    \item Regressores externos: taxa SELIC, câmbio, índice de
    confiança da indústria (FGV).
    \item Otimizar hiperparâmetros via validação cruzada
    (grid search em $\tau$ e $\sigma_s$).
    \item Modelo híbrido: Prophet + ARIMA nos resíduos, para
    capturar a autocorrelação residual detectada.
    \item Tendência logística com capacidade de saturação.
\end{enumerate}

% ============================================================
\section{Conclusão}

O Prophet decompõe séries temporais em tendência (changepoints
bayesianos), sazonalidade (Fourier) e feriados, estimados via MAP.
Na ilustração com a produção industrial brasileira (IBGE, Série
21859, jan/2002 a jan/2026, 289 observações), o modelo capturou
os padrões sazonais e as mudanças de tendência, obtendo MAPE de
3.79\% no teste e 4.88\% na validação cruzada, com cobertura de
100\% do intervalo de confiança de 95\%.

Contudo, os diagnósticos revelaram limitações importantes: embora
os resíduos sigam distribuição normal (Jarque-Bera, $p = 0{,}70$),
o teste $t$ indicou viés sistemático (média dos resíduos $= 3{,}47$,
$p < 0{,}0001$), o teste de Ljung-Box detectou autocorrelação
residual significativa ($p < 0{,}0001$ em todos os lags), e o teste
de Spearman apontou heterocedasticidade ($\rho = 0{,}69$,
$p < 0{,}0001$). O $R^2 = 0{,}57$ no teste indica que o modelo
explica pouco mais da metade da variabilidade, e o Sazonal-Naïve
superou o Prophet em todas as métricas (MAPE de 1.76\% vs.\ 3.79\%).

Esses resultados confirmam que a principal fraqueza do Prophet é
\textbf{não modelar a autocorrelação residual}, tratando
$\varepsilon_t$ como i.i.d. O viés positivo sugere subestimação
sistemática. Melhorias possíveis incluem: modelo híbrido
Prophet+ARIMA nos resíduos, inclusão de COVID-19 como evento
especial, e regressores externos (SELIC, câmbio). Ainda assim,
o modelo se mostrou superior aos benchmarks Naïve, Média e Drift,
e sua força permanece na \textbf{interpretabilidade} dos
componentes e na facilidade de uso.

% ============================================================
\newpage
\section*{Referências}
\addcontentsline{toc}{section}{Referências}

\begin{thebibliography}{99}

% ============================================================
% REFERÊNCIAS DO PROFESSOR (Bibliografia da disciplina)
% ============================================================

\bibitem{morettin2017}
MORETTIN, P.~A.
\textit{Econometria Financeira: Um Curso em Séries Temporais
Financeiras}. 3.ª~ed. Blucher, São Paulo, 2017.
ISBN: 978-85-212-1157-0.

\bibitem{morettin2006}
MORETTIN, P.~A.; TOLOI, C.~M.~C.
\textit{Análise de Séries Temporais}. 2.ª~ed.
ABE --- Projeto Fisher. Blucher, São Paulo, 2006.
ISBN: 978-85-212-0389-6.

\bibitem{morettin2020}
MORETTIN, P.~A.; TOLOI, C.~M.~C.
\textit{Análise de Séries Temporais, Volume 2: Modelos
Multivariados e Não-Lineares}. ABE --- Projeto Fisher. Blucher,
São Paulo, 2020.
ISBN: 978-85-212-1626-1.


% ============================================================
% REFERÊNCIAS DO MODELO PROPHET E COMPLEMENTARES
% ============================================================

\bibitem{taylor2018}
TAYLOR, S.~J.; LETHAM, B.
Forecasting at scale.
\textit{The American Statistician}, v.~72, n.~1, p.~37--45, 2018.
ISSN: 0003-1305.

\bibitem{box2015}
BOX, G.~E.~P.; JENKINS, G.~M.; REINSEL, G.~C.; LJUNG, G.~M.
\textit{Time Series Analysis: Forecasting and Control}.
5th~ed. Wiley, 2015.
ISBN: 978-1-118-67502-1.

\bibitem{harvey1993}
HARVEY, A.~C.; SHEPHARD, N.
Structural time series models. In: \textit{Handbook of
Statistics}, v.~11. Elsevier, 1993.
ISBN: 978-0-444-89572-2.

\bibitem{carpenter2017}
CARPENTER, B. et al.
Stan: A probabilistic programming language.
\textit{Journal of Statistical Software}, v.~76, n.~1, 2017.
ISSN: 1548-7660.

\bibitem{hyndman2021}
HYNDMAN, R.~J.; ATHANASOPOULOS, G.
\textit{Forecasting: Principles and Practice}. 3rd~ed. OTexts,
2021.
ISBN: 978-0-9875071-3-6.

\bibitem{lewis1982}
LEWIS, C.~D.
\textit{Industrial and Business Forecasting Methods}.
Butterworth-Heinemann, 1982.
ISBN: 978-0-408-00532-2.

\bibitem{ibge2025}
IBGE. Pesquisa Industrial Mensal --- Produção Física (PIM-PF).
Disponível em:
\nolinkurl{https://www.ibge.gov.br/estatisticas/economicas/industria/9294-pesquisa-industrial-mensal-producao-fisica-brasil.html}.
Acesso em: jul.\ 2025.

\end{thebibliography}
\end{document}
